% Part: turing-machines
% Chapter: undecidability
% Section: enumerating-tms

\documentclass[../../../include/open-logic-section]{subfiles}

\begin{document}

\olfileid{tur}{und}{enu}
\olsection{تعداد آلات تورنغ}

\begin{explain}
يمكننا أن نبين أن مجموعة جميع آلات تورنغ !!{enumerable}. وينتج ذلك من
إمكان وصف كل آلة تورنغ وصفًا منتهيًا. فمجموعة الحالات وأبجدية الشريط
مجموعتان منتهيتان. ودالة الانتقال دالة جزئية من $Q \times \Sigma$ إلى
$Q \times \Sigma \times \{\TMleft,\TMright,\TMstay\}$، ولذلك يمكن
تحديدها هي الأخرى بإدراج قيمها عند أزواج الوسيطات المنتهية العدد التي
تكون معرّفة عندها.

هذا صحيح في حدوده، لكن ثمة فرق دقيق. فلم يفرض تعريف آلات تورنغ أي
قيد على ال!!{element}s التي يمكن أن تحتويها مجموعة الحالات وأبجدية
الشريط. ولذلك توجد، من الناحية التقنية، آلة تورنغ تستعمل كل عدد حقيقي
معين حالةً من حالاتها. لكن \emph{سلوك} آلة تورنغ مستقل عن الأشياء التي
تؤدي دور الحالات ورموز الأبجدية. انظر إلى آلتي تورنغ في
\olref{fig:variants}.
\begin{figure}
\begin{center}
\begin{tikzpicture}[->,>=stealth',shorten >=1pt,auto,node distance=2.8cm,
                    semithick]
  \tikzstyle{every state}=[fill=none,draw=black,text=black]

  \node[initial,state]         (A)              {$q_0$};
  \node[state]         (B) [right of=A] {$q_1$};

  \path (A) edge [bend left] node {\TMtrans{\TMstroke}{\TMstroke}{\TMright}} (B)
        (B) edge [loop above] node {\TMtrans{\TMblank}{\TMblank}{\TMright}} (B)
            edge [bend left] node {\TMtrans{\TMstroke}{\TMstroke}{\TMright}} (A);
\end{tikzpicture}\\
\begin{tikzpicture}[->,>=stealth',shorten >=1pt,auto,node distance=2.8cm,
  semithick]
\tikzstyle{every state}=[fill=none,draw=black,text=black]

\node[initial,state]         (A)              {$s$};
\node[state]         (B) [right of=A] {$h$};

\path (A) edge [bend left] node {\TMtrans{A}{A}{\TMright}} (B)
(B) edge [loop above] node {\TMtrans{\TMblank}{\TMblank}{\TMright}} (B)
edge [bend left] node {\TMtrans{A}{A}{\TMright}} (A);
\end{tikzpicture}
\end{center}
\caption{صورتان مختلفتان لآلة \emph{الزوجية}}
\ollabel{fig:variants}
\end{figure}
يقابل هذان المخططان آلتين: الآلة $M$، ذات أبجدية الشريط
$\Sigma=\{\TMendtape,\TMblank,\TMstroke\}$ ومجموعة الحالات
$\{q_0,q_1\}$، والآلة $M'$، ذات الأبجدية
$\Sigma'=\{\TMendtape,\TMblank,A\}$ والحالتين $\{s,h\}$. لكن
تعليماتهما متماثلة فيما عدا ذلك: تتوقف $M$ عند متتالية مؤلفة من $n$
من علامات~$\TMstroke$ إذا وفقط إذا كان $n$ زوجيًا، وتتوقف $M'$ عند
متتالية مؤلفة من $n$ من رموز~$A$ إذا وفقط إذا كان $n$ زوجيًا. وكل ما
فعلناه هو إعادة تسمية $\TMstroke$ بـ~$A$، و$q_0$ بـ~$s$، و$q_1$
بـ~$h$. ويمكن تعميم هذا المثال: نستطيع عد آلتي تورنغ متماثلتين ما دامت
إحداهما تنتج من الأخرى بإعادة تسمية الرموز والحالات على هذا النحو.
بل يمكننا ببساطة عد رموز آلة تورنغ وحالاتها أعدادًا صحيحة موجبة: فنفكر
في~$1$ بدلًا من~$\sigma_0$، وفي~$2$ بدلًا من~$\sigma_1$، وهكذا؛
ويكون $\TMendtape$ هو~$1$، و$\TMblank$ هو~$2$، وهكذا. وبهذه الطريقة
تصير آلة \emph{الزوجية} الآلة المرسومة في
\olref{fig:standard-even}.
\begin{figure}
\[\begin{tikzpicture}[->,>=stealth',shorten >=1pt,auto,node distance=2.8cm,
  semithick]
\tikzstyle{every state}=[fill=none,draw=black,text=black]

\node[initial,state]         (A)              {$1$};
\node[state]         (B) [right of=A] {$2$};

\path (A) edge [bend left] node {\TMtrans{3}{3}{\TMright}} (B)
(B) edge [loop above] node {\TMtrans{2}{2}{\TMright}} (B)
edge [bend left] node {\TMtrans{3}{3}{\TMright}} (A);
\end{tikzpicture}
\]
\caption{آلة \emph{زوجية} قياسية}
\ollabel{fig:standard-even}
\end{figure}
يمكن أن نسمي آلة تورنغ التي تكون حالاتها ورموزها أعدادًا صحيحة موجبة
آلة \emph{قياسية}، وأن نقتصر من الآن فصاعدًا على الآلات القياسية.
\footnote{مصطلح «الآلة القياسية» ليس مصطلحًا قياسيًا.}

أردنا أن نبين أن مجموعة آلات تورنغ !!{enumerable}، وفي ضوء الاعتبارات
السابقة يكفي أن نبين أن مجموعة آلات تورنغ القياسية !!{enumerable}.
لنفترض أننا أعطينا آلة تورنغ قياسية
$M=\tuple{Q,\Sigma,q_0,\delta}$. كيف يمكننا وصفها بسلسلة منتهية من
الأعداد الصحيحة الموجبة؟ ندرج أولًا عدد الحالات، ثم الحالات نفسها، ثم
عدد الرموز، ثم الرموز نفسها، ثم الحالة الابتدائية. (تذكر أن هذه كلها
أعداد صحيحة موجبة، لأن $M$ آلة قياسية.) وماذا عن~$\delta$؟ إن مجموعة
الوسيطات الممكنة، أي الأزواج $\tuple{q,\sigma}$، منتهية لأن $Q$
و~$\Sigma$ منتهيتان. ولذلك لا تتعدى المعلومات التي تحملها~$\delta$
القائمة المنتهية لجميع الخماسيات $\tuple{q,\sigma,q',\sigma',d}$ التي
تكون عندها
$\delta(q,\sigma)=\tuple{q',\sigma',D}$، حيث $d$ عدد يرمز الاتجاه~$D$
(وليكن $1$ لـ~$\TMleft$، و$2$ لـ~$\TMright$، و$3$ لـ~$\TMstay$).

وبهذه الطريقة يمكن وصف كل آلة تورنغ قياسية بقائمة منتهية من الأعداد
الصحيحة الموجبة، أي بمتتالية $s_M\in(\PosInt)^*$. فمثلًا ترمز آلة
\emph{الزوجية} القياسية بالمتتالية
\[
2, \underbrace{1, 2}_Q, 3, \overbrace{1, 2, 3}^\Sigma, 1, \underbrace{1, 3, 2, 3, 2}_{\delta(1,3) = \tuple{2,3,R}},
\overbrace{2, 2, 2, 2, 2}^{\delta(2,2) = \tuple{2,2,R}},
\underbrace{2, 3, 1, 3, 2}_{\delta(2,3) = \tuple{1,3,R}}.
\]
\end{explain}

\begin{thm}
توجد دوال من $\Nat$ إلى~$\Nat$ غير قابلة للحساب بآلة تورنغ.
\end{thm}

\begin{proof}
نعلم أن مجموعة المتتاليات المنتهية من الأعداد الصحيحة الموجبة
$(\PosInt)^*$ هي !!{enumerable}
(\cref{sfr:siz:zigzag:prob:posint-star}). وينتج من ذلك أن مجموعة أوصاف
آلات تورنغ القياسية، بوصفها مجموعة جزئية من~$(\PosInt)^*$، قابلة
للتعداد هي أيضًا. وكل دالة من $\Nat$ إلى~$\Nat$ قابلة للحساب بآلة
تورنغ تحسبها آلة تورنغ ما (بل آلات كثيرة). وبإعادة تسمية حالاتها
ورموزها بأعداد صحيحة موجبة (وبوجه خاص تسمية $\TMendtape$ بالعدد~$1$،
و$\TMblank$ بالعدد~$2$، و$\TMstroke$ بالعدد~$3$) نرى أن آلة تورنغ
قياسية تحسب كل دالة قابلة للحساب بآلة تورنغ. وهذا يعني أن مجموعة جميع
الدوال من $\Nat$ إلى~$\Nat$ القابلة للحساب بآلة تورنغ قابلة للتعداد
أيضًا.

ومن ناحية أخرى، ليست مجموعة جميع الدوال من $\Nat$ إلى~$\Nat$
!!{enumerable} (\cref{sfr:siz:red:prob:nat-nat}). فلو كانت كل الدوال
قابلة للحساب بآلة تورنغ ما، لأمكننا تعداد مجموعة جميع الدوال بإدراج
أوصاف آلات تورنغ التي تحسبها. ومن ثم توجد دوال غير قابلة للحساب بآلة
تورنغ.
\end{proof}

\begin{prob}
  هل يمكنك أن تجد طريقة لوصف آلات تورنغ لا تتطلب إدراج الحالات ورموز
  الأبجدية صراحة؟ يجوز لك أن تعرّف مفهومك الخاص للآلة «القياسية»، لكن
  وضح لماذا تستطيع آلة «قياسية» بالمعنى الجديد أن تحاكي كل آلة تورنغ.
\end{prob}

\end{document}
